% sections/zh/5-experiments.tex -- 按规格 6C 的结构化骨架（5.1–5.7 + 算例因子效应）；数值全部为占位符。
\section{实验}\label{sec:experiments}

\draftnote{本节全部数值为占位符，将由 \texttt{paper\_assets/scripts/make\_tables.py} 从 \texttt{result/*.csv} 按《实验设计规格》6B 契约填充；全部图形为占位框，由 opt-paper-scifig 生成后替换。方向性表述在裁决完成前均为预期。}

\subsection{实验设置}\label{sec:settings}
评测集是由算例参数表生成的固定设计集：平均多重度 $\bar N$ 为规模因子，多重度离散度 $V_N$、柔性度 $F$ 与工时异质度 $V_p$ 为三个结构因子，类型数 $R = 8$、机器数 $m = 10$（表 \ref{tab:instances}）。主实验档为 $\bar N = 10$ 下的 $3\times3\times3$ 全因子；small 档由 CP-SAT\cite{ortools} 在 1800 s 内给出精确参考解；large 档把多重度外推到训练区间之外；Brandimarte 算例\cite{brandimarte1993routing} 为 $N_r = 1$ 的特例。算例以文件形式固定并随论文发布，复现依靠这些文件与 \NumRuns{} 次独立运行而非随机种子。表 \ref{tab:hyper} 列出超参数，表 \ref{tab:baselines} 列出基线：每族一个强代表，取推荐配置与同等墙钟预算。指标为 makespan、相对精确解或参考解的最优性 gap、用于跨档比较的流体 gap $(C_{\max} - L^F)/L^F$、可行率、1 / 10 / 60 s 检查点下的质量-时间曲线面积，以及精确解码次数与达到目标的时间。显著性用逐算例 run 均值的配对 Wilcoxon 检验（$p < 0.05$）评估，方法数不少于四时另给 Friedman–Nemenyi 临界差图。

\begin{table}[htbp]
\centering\small
\caption{评测算例设置}\label{tab:instances}
% generated by paper_assets/scripts/make_tables.py -- do not edit by hand

\begin{tabularx}{\linewidth}{@{}l>{\hsize=1.000\hsize\raggedright\arraybackslash}X>{\hsize=0.900\hsize\raggedright\arraybackslash}Xl>{\hsize=1.100\hsize\raggedright\arraybackslash}X@{}}
\toprule
\textbf{\bi{档位}{Tier}} & \textbf{\bi{规模因子 $\bar N$}{Scale factor $\bar N$}} & \textbf{\bi{结构因子}{Structural factors}} & \textbf{\bi{算例数}{Instances}} & \textbf{\bi{用途}{Purpose}} \\
\midrule
Small & 2, 4 & \bi{各取低 / 高两端}{low / high ends} & 16 & \bi{精确参考解校准真实 gap}{exact reference solutions calibrate the true gap} \\
Main & 10 & \bi{3 $\times$ 3 $\times$ 3 全因子}{$3\times3\times3$ full factorial} & 27 & \bi{主对比 + 因子效应分析}{main comparison and factor-effect analysis} \\
Large & 25, 50 & \bi{各取低 / 高两端}{low / high ends} & 16 & \bi{多重度外推（商图表示直接求解）}{multiplicity extrapolation (quotient representation solved directly)} \\
Benchmark & \bi{Brandimarte mk01–mk10（$N_r = 1$ 退化为标准 FJSP）}{Brandimarte mk01--mk10 ($N_r = 1$, standard FJSP)} & \bi{外部给定}{given} & 10 & \bi{公开基准（补充，不替代设计集）}{public benchmark (supplementary to the design set)} \\
\bottomrule
\end{tabularx}

\end{table}

\begin{table}[htbp]
\centering\small
\caption{超参数设置（草案值；终值以确认后的训练参数表为准）}\label{tab:hyper}
% generated by paper_assets/scripts/make_tables.py -- do not edit by hand

\begin{tabularx}{\linewidth}{@{}l>{\hsize=0.850\hsize\raggedright\arraybackslash}X>{\hsize=1.150\hsize\raggedright\arraybackslash}X@{}}
\toprule
\textbf{\bi{类别}{Group}} & \textbf{\bi{参数}{Parameter}} & \textbf{\bi{取值}{Value}} \\
\midrule
\bi{训练}{Training} & \texttt{train.paradigm} & \bi{generalist（排序器离线训练、冻结后零样本）}{generalist (ranker trained offline, frozen, zero-shot)} \\
\bi{训练}{Training} & \texttt{instances\_\allowbreak{}per\_\allowbreak{}epoch} & \bi{256 个随机算例，FRES-lite 自举 32\,768 个候选-标签对}{256 random instances; 32\,768 candidate--label pairs bootstrapped by FRES-lite} \\
\bi{训练}{Training} & \texttt{epochs} /\allowbreak{} \texttt{plateau\_\allowbreak{}patience} & 60 / 10 \\
\bi{训练}{Training} & \texttt{wall\_\allowbreak{}budget\_\allowbreak{}s} /\allowbreak{} \texttt{eval\_\allowbreak{}every} /\allowbreak{} \texttt{env\_\allowbreak{}batch} & \bi{待实测}{to be measured} \\
\bi{算法}{Optimizer} & \texttt{lr} & \bi{1e-3，余弦退火}{1e-3, cosine annealing} \\
\bi{算法}{Optimizer} & \texttt{batch\_\allowbreak{}size} /\allowbreak{} \texttt{loss} & \bi{512 / pairwise ranking（logistic）}{512 / pairwise ranking (logistic)} \\
\bi{网络}{Network} & \texttt{embed\_\allowbreak{}dim} /\allowbreak{} \texttt{n\_\allowbreak{}layers} & \bi{128 / 3（异构消息传递）}{128 / 3 (heterogeneous message passing)} \\
\bi{宿主}{Host EA} & \texttt{population} /\allowbreak{} \texttt{offspring} & 100 / 100 \\
\bi{宿主}{Host EA} & $B$ / $K_L$ / $K_H$ / $K_R$ & 96 / 10 / 5 / 3 \\
\bi{宿主}{Host EA} & \bi{配额迁移步长 $b$ / 优先键扰动幅度}{quota step $b$ / priority-key perturbation} & 1 / 0.1 \\
\bottomrule
\end{tabularx}

\end{table}

\begin{table}[htbp]
\centering\small
\caption{对比方法及配置}\label{tab:baselines}
% generated by paper_assets/scripts/make_tables.py -- do not edit by hand

\begin{tabularx}{\linewidth}{@{}>{\hsize=0.700\hsize\raggedright\arraybackslash}Xl>{\hsize=1.300\hsize\raggedright\arraybackslash}X@{}}
\toprule
\textbf{\bi{方法}{Method}} & \textbf{\bi{类别}{Class}} & \textbf{\bi{实现与配置}{Implementation and configuration}} \\
\midrule
\bi{HGA（job-indexed）}{HGA (job-indexed)} & \bi{元启发式}{metaheuristic} & \bi{文献推荐参数（双层编码遗传算法）；种群 100、POX 交叉 0.8、变异 0.15、精英 10\%；同墙钟；每算例 3 次重启}{recommended settings of the two-layer encoded GA; population 100, POX crossover 0.8, mutation 0.15, elitism 10\%; same wall clock; 3 restarts per instance} \\
\bi{关键块禁忌搜索（TS）}{Critical-block tabu search (TS)} & \bi{元启发式}{metaheuristic} & \bi{Mastrolilli–Gambardella 型，自实现并在 mk 基准核对；禁忌长度 12、关键块邻域 + 关键工序改派；同墙钟；3 次重启}{Mastrolilli--Gambardella type, reimplemented and checked on the mk benchmark; tenure 12, critical-block neighbourhood with reassignment; same wall clock; 3 restarts} \\
\bi{流体进化算法（FMAE-s）}{Fluid-assisted EA (FMAE-s)} & \bi{元启发式}{metaheuristic} & \bi{流体主从进化算法的无换模静态版（复现）；原文推荐参数；流体比例只作初始化与规则；同墙钟；3 次重启}{static, setup-free version of the fluid master--slave EA (reimplementation); published settings; fluid ratios used only for initialization and rules; same wall clock; 3 restarts} \\
DSLFMAE & \bi{学习型}{learning-based} & \bi{PPO 辅助流体进化算法（复现或机制等价复现，标注 reimplementation）；原文推荐配置；与本方法同等训练预算，3 run}{PPO-assisted fluid EA (reimplementation, marked as such); published configuration; same training budget as ours, 3 runs} \\
HGNN-DRL & \bi{学习型}{learning-based} & \bi{官方代码（异构图神经网络 + PPO 构造式调度）；推荐配置；贪心 + 采样两口径；同等训练预算，3 run}{official code (heterogeneous GNN with PPO, constructive); recommended configuration; greedy and sampling decoding; same training budget, 3 runs} \\
CP-SAT & \bi{参考解}{reference} & \bi{OR-Tools；1800 s（small）/ 与元启发式同墙钟（main，anytime 参考）}{OR-Tools; 1800 s (small) / same wall clock as the metaheuristics (main, anytime reference)} \\
\bottomrule
\end{tabularx}

\end{table}

\subsection{与基线方法的对比}\label{sec:comparison}
本小节回答 FRES 是否优于元启发式与学习型基线、差异是否显著。表 \ref{tab:main} 报告 small 与 main 档在 10 s 与 60 s 检查点的结果；图 \ref{fig:main} 给出 gap 分布与临界差图。\pending{仅 FRES-lite 就达到不高于 HGA 与 FMAE-s 的 gap，完整方法在主实验档相对最强基线改善 \ImpvOverBest{}。}

\begin{table}[htbp]
\centering\small
\caption{主要结果：\NumRuns{} 次运行的均值 $\pm$ 标准差；最优加粗；$\dagger$ 表示显著优于次优（Wilcoxon，$p < 0.05$）。S = small 档（相对精确解的 gap），M = main 档（相对参考解的 gap）}\label{tab:main}
% generated by paper_assets/scripts/make_tables.py -- do not edit by hand

\begin{tabularx}{\linewidth}{@{}>{\hsize=1.500\hsize\raggedright\arraybackslash}X>{\hsize=0.800\hsize\raggedright\arraybackslash}X>{\hsize=0.800\hsize\raggedright\arraybackslash}X>{\hsize=1.000\hsize\raggedright\arraybackslash}X>{\hsize=1.000\hsize\raggedright\arraybackslash}X>{\hsize=0.800\hsize\raggedright\arraybackslash}X>{\hsize=1.100\hsize\raggedright\arraybackslash}X@{}}
\toprule
\textbf{\bi{方法}{Method}} & \textbf{\bi{S-gap（\%）}{S gap (\%)}} & \textbf{\bi{S-可行率}{S feas.}} & \textbf{\bi{M-gap 10 s（\%）}{M gap 10 s (\%)}} & \textbf{\bi{M-gap 60 s（\%）}{M gap 60 s (\%)}} & \textbf{\bi{M-可行率}{M feas.}} & \textbf{\bi{解码次数（$\times10^{3}$）}{Decodings ($\times10^{3}$)}} \\
\midrule
\bi{HGA（job-indexed）}{HGA (job-indexed)} & \phpm & \ph & \phpm & \phpm & \ph & \ph \\
\bi{关键块禁忌搜索（TS）}{Critical-block tabu search (TS)} & \phpm & \ph & \phpm & \phpm & \ph & \ph \\
\bi{流体进化算法（FMAE-s）}{Fluid-assisted EA (FMAE-s)} & \phpm & \ph & \phpm & \phpm & \ph & \ph \\
\bi{DSLFMAE}{DSLFMAE} & \phpm & \ph & \phpm & \phpm & \ph & \ph \\
\bi{HGNN-DRL}{HGNN-DRL} & \phpm & \ph & \phpm & \phpm & \ph & \ph \\
\bi{CP-SAT（同墙钟）}{CP-SAT (same wall clock)} & \phpm & \ph & \phpm & \phpm & \ph & \ph \\
\bi{\textbf{FRES}}{\textbf{FRES}} & \phb & \phb & \phb & \phb & \phb & \phb \\
\bottomrule
\end{tabularx}

\end{table}

\begin{figure}[htbp]
\figplace{主实验档 60 s 时各方法的 gap 箱线图，每档一个面板，带 Wilcoxon 显著性标注；内嵌临界差图}{\texttt{result/*/eval\_results.csv}（gap、method、tier）；$p$ 值与秩取自 \texttt{result/stats\_summary.csv}}{opt-paper-scifig 路线 A（matplotlib），栏宽 88\,mm}{52mm}
\caption{主实验档的 gap 分布与临界差图}\label{fig:main}
\end{figure}

\subsection{消融与排序研究}\label{sec:ablation}
本小节回答增益是否来自三个机制、排序器是否减少了无效解码。变体 A1--A5 逐一去掉流体配额、配额压缩、排序器、残差限定与保底（表 \ref{tab:ablation}）。图 \ref{fig:ablation} 给出各变体的收敛过程，图 \ref{fig:recall} 给出同候选预算下排序器相对流体启发式与通用排序器的 top-10 recall，图 \ref{fig:safeguard} 给出被接受子代来自并集各源的占比。\pending{完整方法比 FRES-lite 少解码 \OursDecodeReduction{} 个候选而质量不降。}

\begin{table}[htbp]
\centering\small
\caption{主实验档的消融与机制指标}\label{tab:ablation}
% generated by paper_assets/scripts/make_tables.py -- do not edit by hand

\begin{tabularx}{\linewidth}{@{}>{\hsize=1.000\hsize\raggedright\arraybackslash}Xllll@{}}
\toprule
\textbf{\bi{变体}{Variant}} & \textbf{\bi{gap（\%）}{gap (\%)}} & \textbf{\bi{可行率}{feas.}} & \textbf{top-10 recall} & \textbf{\bi{解码次数（$\times10^{3}$）}{Decodings ($\times10^{3}$)}} \\
\midrule
\bi{完整方法（FRES）}{Full method (FRES)} & \phb & \phb & \phb & \phb \\
\bi{A1 无流体配额}{A1 no fluid quota} & \phpm & \ph & \ph & \ph \\
\bi{A2 无配额压缩}{A2 no quota compression} & \phpm & \ph & \ph & \ph \\
\bi{A3 无排序器（FRES-lite）}{A3 no ranker (FRES-lite)} & \phpm & \ph & \ph & \ph \\
\bi{A4 通用排序器}{A4 generic ranker} & \phpm & \ph & \ph & \ph \\
\bi{A5 无保底}{A5 no safeguard} & \phpm & \ph & \ph & \ph \\
\bottomrule
\end{tabularx}

\end{table}

\begin{figure}[htbp]
\figplace{FRES 与 A1--A5 的在位最优 gap 随墙钟变化，均值线加 $\pm$ 标准差阴影带，时间轴取对数}{\texttt{result/*/curve.csv} 跨 run 聚合}{opt-paper-scifig 路线 A，栏宽 88\,mm}{50mm}
\caption{消融变体的收敛过程}\label{fig:ablation}
\end{figure}

\begin{figure}[htbp]
\figplace{真实改善动作的 top-10 recall 随解码候选数变化，三条曲线（残差排序器、流体启发式、通用排序器），同一候选池}{\texttt{result/*/log.csv} 的 ranking 指标与 \texttt{operator\_trace.csv}}{opt-paper-scifig 路线 A，栏宽 88\,mm}{50mm}
\caption{同候选预算下的排序质量}\label{fig:recall}
\end{figure}

\begin{figure}[htbp]
\figplace{被接受子代的来源（学习器 / 流体启发式 / 随机）随代数变化的堆叠柱，每个 $V_N$ 水平一个面板}{\texttt{operator\_trace.csv}（best\_from、generation）join \texttt{index.csv}}{opt-paper-scifig 路线 A，栏宽 88\,mm}{48mm}
\caption{保底并集下被接受子代的来源}\label{fig:safeguard}
\end{figure}

\subsection{参数敏感性}\label{sec:sensitivity}
本小节回答预选参数是否需要调参。图 \ref{fig:sens} 在验证档上扫描候选池大小 $B$ 与学习器份额 $K_L$。\pending{性能在扫描网格上平坦。}

\begin{figure}[htbp]
\figplace{主实验档 gap 在 $B \in \{48, 96, 192\}$ 与 $K_L \in \{8, 16, 24\}$ 上的热力图，格内标均值，选定设置加框}{\texttt{result/sensitivity.csv}}{opt-paper-scifig 路线 A，栏宽 88\,mm}{48mm}
\caption{对候选池大小与学习器份额的敏感性}\label{fig:sens}
\end{figure}

\subsection{算例因子效应}\label{sec:factor}
本小节回答方法在什么问题特性下更强或更弱、理论界是否得到支持。它与上一小节的参数敏感性不同：这里的因子属于问题而非算法。图 \ref{fig:rounding} 在 $\bar N$ 与 $V_p$ 上比较实证取整差距 $G_Q/L^F$ 与命题 \ref{prop:rounding} 的界 $\bar B/L^F$（猜想 \ref{conj:rounding}）；表 \ref{tab:factor} 与图 \ref{fig:factor} 给出相对最强基线增益的主效应与双因子交互（猜想 \ref{conj:gain}）。误差棒为运行间标准差，不是单元内重复。

\begin{table}[htbp]
\centering\small
\caption{主实验档的算例因子主效应（27 个参数各异的算例，每水平 9 个；误差为 \NumRuns{} 次运行的运行间标准差）}\label{tab:factor}
% generated by paper_assets/scripts/make_tables.py -- do not edit by hand

\begin{tabularx}{\linewidth}{@{}>{\hsize=0.900\hsize\raggedright\arraybackslash}Xll>{\hsize=1.100\hsize\raggedright\arraybackslash}Xl@{}}
\toprule
\textbf{\bi{因子}{Factor}} & \textbf{\bi{水平}{Level}} & \textbf{\bi{gap（\%）}{gap (\%)}} & \textbf{\bi{相对最强基线的增益（pp）}{Gain over the strongest baseline (pp)}} & \textbf{\bi{解码次数（$\times10^{3}$）}{Decodings ($\times10^{3}$)}} \\
\midrule
\bi{多重度离散度 $V_N$}{Multiplicity dispersion $V_N$} & 0.0 & \phpm & \ph & \ph \\
\bi{多重度离散度 $V_N$}{Multiplicity dispersion $V_N$} & 0.5 & \phpm & \ph & \ph \\
\bi{多重度离散度 $V_N$}{Multiplicity dispersion $V_N$} & 1.0 & \phpm & \ph & \ph \\
\bi{柔性度 $F$}{Flexibility $F$} & 0.3 & \phpm & \ph & \ph \\
\bi{柔性度 $F$}{Flexibility $F$} & 0.55 & \phpm & \ph & \ph \\
\bi{柔性度 $F$}{Flexibility $F$} & 0.8 & \phpm & \ph & \ph \\
\bi{工时异质度 $V_p$}{Processing-time heterogeneity $V_p$} & 0.2 & \phpm & \ph & \ph \\
\bi{工时异质度 $V_p$}{Processing-time heterogeneity $V_p$} & 0.5 & \phpm & \ph & \ph \\
\bi{工时异质度 $V_p$}{Processing-time heterogeneity $V_p$} & 0.8 & \phpm & \ph & \ph \\
\bottomrule
\end{tabularx}

\end{table}

\begin{figure}[htbp]
\figplace{两个面板：实证 $G_Q/L^F$ 与界 $\bar B/L^F$ 随 $\bar N$（对数轴，small / main / large 档）与随 $V_p$（主实验档）的变化，标记点带运行间误差棒}{\texttt{relaxation\_diag.csv}（rounding\_gap、rounding\_bound、fluid\_lb）join \texttt{index.csv}}{opt-paper-scifig 路线 A，栏宽 88\,mm}{48mm}
\caption{实证取整差距与其理论界}\label{fig:rounding}
\end{figure}

\begin{figure}[htbp]
\figplace{左：相对最强基线增益在 $V_N$、$F$、$V_p$ 上的主效应图，带误差棒；右：双因子交互热力图}{\texttt{eval\_results.csv} 按 instance\_id join \texttt{index.csv}}{opt-paper-scifig 路线 A，双栏宽 178\,mm}{52mm}
\caption{相对最强基线增益的因子效应}\label{fig:factor}
\end{figure}

\subsection{训练与搜索行为}\label{sec:training}
本小节回答搜索与排序器训练是否稳定。图 \ref{fig:anytime} 给出固定验证集上各方法在位最优 gap 随墙钟的变化与排序器的训练曲线。

\begin{figure}[htbp]
\figplace{左：各方法验证集 gap 随墙钟的 anytime 曲线，均值 $\pm$ 标准差；右：排序器 \NumRuns{} 次训练的逐 epoch 排序损失与 top-10 recall}{\texttt{curve.csv} 与 \texttt{log.csv}}{opt-paper-scifig 路线 A，栏宽 88\,mm}{50mm}
\caption{anytime 行为与排序器训练曲线}\label{fig:anytime}
\end{figure}

\subsection{向更高多重度的泛化}\label{sec:generalization}
本小节回答方法是否只在训练多重度区间内有效。表 \ref{tab:generalization} 用冻结的排序器直接求解 large 档（$\bar N \in \{25, 50\}$）；图 \ref{fig:gen} 画出 gap 随 $\bar N$ 的变化并以阴影标出训练区间。\pending{FRES 的 gap 在 $\bar N$ 上平坦，而 job-indexed 宿主的 gap 随之增长。}

\begin{table}[htbp]
\centering\small
\caption{训练区间（$\bar N \le 12$）之外的多重度外推}\label{tab:generalization}
% generated by paper_assets/scripts/make_tables.py -- do not edit by hand

\begin{tabularx}{\linewidth}{@{}>{\hsize=1.100\hsize\raggedright\arraybackslash}X>{\hsize=1.000\hsize\raggedright\arraybackslash}X>{\hsize=0.900\hsize\raggedright\arraybackslash}X>{\hsize=1.000\hsize\raggedright\arraybackslash}X>{\hsize=0.900\hsize\raggedright\arraybackslash}X@{}}
\toprule
\textbf{\bi{方法}{Method}} & \textbf{\bi{$\bar N = 25$ gap（\%）}{$\bar N = 25$ gap (\%)}} & \textbf{\bi{$\bar N = 25$ 可行率}{$\bar N = 25$ feas.}} & \textbf{\bi{$\bar N = 50$ gap（\%）}{$\bar N = 50$ gap (\%)}} & \textbf{\bi{$\bar N = 50$ 可行率}{$\bar N = 50$ feas.}} \\
\midrule
\bi{HGA（job-indexed）}{HGA (job-indexed)} & \phpm & \ph & \phpm & \ph \\
\bi{流体进化算法（FMAE-s）}{Fluid-assisted EA (FMAE-s)} & \phpm & \ph & \phpm & \ph \\
\bi{A2 无配额压缩}{A2 no quota compression} & \phpm & \ph & \phpm & \ph \\
\bi{\textbf{FRES}}{\textbf{FRES}} & \phb & \phb & \phb & \phb \\
\bottomrule
\end{tabularx}

\end{table}

\begin{figure}[htbp]
\figplace{FRES、HGA、FMAE-s 与 A2 的 gap 随 $\bar N$ 变化，训练区间以阴影标出，$x$ 轴取对数}{\texttt{eval\_results.csv}（tier、gap）join \texttt{index.csv}（scale）}{opt-paper-scifig 路线 A，栏宽 88\,mm}{48mm}
\caption{gap 随平均多重度的变化}\label{fig:gen}
\end{figure}

\subsection{案例分析}\label{sec:case}
本小节回答调度是否合理、配额是否贴近流体分流。图 \ref{fig:case} 给出一个主实验档算例的甘特图，并列出每个类型工序的整数配额与流体比例。

\begin{figure}[htbp]
\figplace{左：FRES 在一个主实验档算例上的甘特图，同类副本同一色系，关键路径加框；右：每个类型工序的 $n_{rki}/N_r$ 与 $q^*_{rki}/N_r$ 对照柱}{单个算例的调度导出；\texttt{relaxation\_diag.csv}}{opt-paper-scifig 路线 A，双栏宽 178\,mm}{60mm}
\caption{单算例的调度与配额-流体对照}\label{fig:case}
\end{figure}
